% !TeX root = ../paper.tex

\section{Introduction}
\label{sec:intro}

Modern Large Language Models (LLMs) can generate source code from
natural-language descriptions and software requirements~\citep{han2024archcode}.
However, functional correctness alone is not sufficient for industrial software.
ISO/IEC~25010 treats functional correctness as only one subcharacteristic within
a broader product quality model used to specify, measure, and evaluate software
quality~\citep{iso25010}. A program can therefore pass functional tests and
still be poorly structured, creating technical debt.

\subsection{Problem Statement}
A prominent line of feedback-based approaches for LLM code generation and repair
relies on execution-derived signals, such as failed tests and runtime
errors, to guide iterative refinement~\citep{chen2024selfdebug}.
A recent multi-dimensional evaluation, by contrast, applies structural quality
indicators such as maintainability and complexity to code that has already been
generated~\citep{li2026multidim},
rather than using them as a feedback signal during generation. This leaves open whether feeding
such structural metrics back into the generation loop can improve the
maintainability of LLM-generated Python code without reducing its functional
correctness.

\subsection{Research Question}
Can iterative static analysis, used as feedback, improve the maintainability of
LLM-generated Python code without reducing its functional correctness?

\noindent The following sub-questions refine this:
\begin{itemize}
  \item What maintainability values do LLM-generated Python solutions reach in a
    baseline condition without feedback?
  \item How do the Maintainability Index~\citep{coleman1994maintainability} and
    complexity metrics change after up
    to three static-analysis feedback iterations?
  \item Does the functional correctness of the solutions change as a result of
    the feedback iterations?
  \item What is the explanatory power of the chosen metrics for this
    investigation?
\end{itemize}

\subsection{Objectives and Hypotheses}
The objective is the design, implementation, and evaluation of a reproducible
experiment for the quality assessment of LLM-generated Python code.
\begin{description}
  \item[H1] Iterative static-analysis feedback increases the aggregate
    Maintainability Index relative to a no-feedback baseline.
  \item[H2] Iterative static-analysis feedback reduces structural complexity,
    measured directly by cyclomatic complexity and cognitive complexity.
  \item[H3] Functional correctness is preserved under the feedback iterations; it
    does not significantly decrease relative to the baseline.
\end{description}

\noindent H1 and H2 are related but tested at different levels: the Maintainability
Index~\citep{coleman1994maintainability} is an aggregate score that already
incorporates complexity and code size, whereas H2 targets the structural
complexity metrics directly---cyclomatic complexity~\citep{mccabe1976complexity}
and cognitive complexity~\citep{cognitivecomplexity}. The individual MI
inputs---Halstead volume, lines of code, and comment density---are examined
descriptively alongside the results (\cref{sec:supp}) rather than promoted to
separate hypotheses, consistent with treating the aggregate Index as a secondary
outcome.

\subsection{Structure of the Paper}
\Cref{sec:related} reviews software-quality metrics and prior work.
\Cref{sec:method} describes the experimental design and feedback pipeline.
\Cref{sec:results} reports and discusses the results by hypothesis.
\Cref{sec:limitations} states the study's limitations and the future work they
suggest. \Cref{sec:conclusion} concludes.
